\section{The small-flow bypass}
\label{sec:bypass}

A \emph{small flow} is one whose size is at most a threshold, 10\,KB in most of our experiments. The sending host marks every packet of a small flow, and its ACKs, with a DSCP code point (AF11), so routers can tell small-flow packets apart without inspecting anything else. Every router forwards a marked packet to a neighbor one hop closer to the destination, hashing the flow over all such neighbors (ECMP over shortest paths), so the packet stays on shortest paths end to end. Between neighboring routers that is the direct link.

Large flows keep using Spraypoint unchanged. The bypass adds no links and does not change Spraypoint's rings or tables; it is a separate lane on the same wires, as a separate VRF or traffic class would be on a real switch.

\paragraph{Why not bypass only neighbors?} A narrower design, which we call the \emph{adjacent bypass}, sends a small flow over the direct link when the two routers are neighbors and sprays it as usual otherwise. It is cheaper on a switch: $d$ extra rules on the server-facing ports and no new table. We implemented it and set it aside, for three reasons. First, its reach is small: only about $d/(n-1)$ of router pairs are neighbors (6.4\% at the paper's scale), so it leaves almost every small flow on a sprayed path. Second, for the pairs it does serve, it gives exactly the same one-hop path as the shortest-path bypass, so it gains nothing over it there. Third, it is not safer: neighbors have a single shortest path ($c = 1$), and the models show that single-path pairs are the ones that overload first (\secref{sec:models}). At the paper's scale, both designs reach their limit at about the same share of small bytes, since the neighbors' direct links set it. The rest of this report uses the shortest-path bypass only.

\subsection{Overhead on a switch (H4)}
\label{sec:h4}

Spraypoint selects its spray or point table by ingress port. The bypass adds a DSCP classification and a shortest-path ECMP table with up to $n-1$ destination entries per router. Its ECMP groups compete for the same resources that constrain Spraypoint's pointing groups (§7.4 of~\cite{rng}). A switch with sufficient table and action capacity could make this decision in one pipeline pass. The number of physical stages and the processing delay depend on the target and its resource occupancy. The simulator models queueing and forwarding paths without switch processing delay, so a target compilation and hardware measurements are still needed to evaluate H4. See \file{docs/OVERHEAD.md} for the pipeline and state analysis.
